% =============================================================================
% MIVA-KNIGHT — Project Proposal (Aligned with Repository Evidence)
% Revises claims in IT571__Project_Proposal (1).pdf to match implemented scope
% documented across capstone .tex sources and companion .ipynb notebooks.
% =============================================================================
\documentclass[12pt,a4paper]{article}

\usepackage[utf8]{inputenc}
\usepackage[T1]{fontenc}
\usepackage[margin=2.5cm]{geometry}
\usepackage{amsmath,amssymb}
\usepackage{booktabs}
\usepackage{enumitem}
\usepackage{hyperref}
\usepackage{xcolor}
\usepackage{setspace}
\setstretch{1.1}

\hypersetup{colorlinks=true,linkcolor=blue!40!black,urlcolor=blue!50!black,
  pdftitle={MIVA-KNIGHT Aligned Project Proposal}}

\title{\textbf{MIVA-KNIGHT}\\[0.25em]
\large Aligned Project Proposal:\\
Multimodal Perception, Hybrid Retrieval, and Grounded Voice QA}
\author{IT Capstone --- revised to match repository artifacts}
\date{\today}

\begin{document}
\maketitle

\begin{abstract}
This proposal supersedes earlier draft language that mixed aspirational benchmarks,
internally inconsistent retrieval weights, and medical datasets not used in the
implemented radiology track. It describes MIVA-KNIGHT as a \emph{research prototype}:
frozen backbone encoders (text, image, audio, and optionally video frames), trainable
projection heads into a shared $512$-dimensional unit hypersphere, FAISS-based dense
retrieval, lightweight knowledge-graph re-ranking from corpus-derived entity
co-occurrence, tiered confidence gating, and LLM generation constrained on retrieved
context. Affect estimation (RAVDESS, CREMA-D) and multimodal sentiment fusion
(CMU-MOSI / Pipeline~E) are \emph{separate training tracks} integrated at the prompt or
routing layer, not a single end-to-end cross-attention model over all corpora.
Companion notebooks exercise representation, fusion, and reinforcement-learning
\emph{demos}; they do not by themselves establish production safety or cloud-free
deployment.
\end{abstract}

\section{Revision note (what changed vs.\ the prior proposal PDF)}

The anchor document (\texttt{IT571\_\_Project\_Proposal (1).pdf}) correctly motivates
hybrid RAG, multimodal inputs, and domain adaptation, but contained several items that
\emph{do not match} the codebase and technical write-ups in this repository:

\begin{itemize}[nosep]
  \item \textbf{Retrieval weights:} objectives stated both $60\%/40\%$ (dense/KG) and
        later $70\%/30\%$. The implementation narrative is better described as
        \emph{two-stage retrieval}: dense shortlist (FAISS) followed by graph-based
        re-scoring, not a single fixed linear blend of uncertain calibration.
  \item \textbf{``Cross-attention fusion'' across COCO + RAVDESS:} the documented system
        primarily uses \emph{modality-specific encoders + shared embedding space} and
        separate affect pipelines. Joint contrastive training of unrelated modalities
        without true pairs was identified as harmful; Pipeline~C is explicitly
        \emph{standalone} audio training. Transformer-style fusion appears for CMU-MOSI
        (Pipeline~E), not as the main COCO/RAVDESS integration path.
  \item \textbf{Medical imaging dataset:} the PDF cites MIMIC-CXR with a
        $\mathrm{P@}1 > 95\%$ goal. The radiology alignment work in-repo centres on
        \textbf{ROCO} (image--caption pairs), with reported retrieval metrics on ROCO
        validation splits --- not MIMIC-CXR training as the primary corpus.
  \item \textbf{External benchmarks:} targets such as HaluBench ($>90\%$) and Diffbot
        multi-hop ($>80\%$) are not evidenced in the repository deliverables and should
        be treated as \textbf{future evaluation}, not capstone commitments.
  \item \textbf{COCO scale:} ``25{,}000 validation images'' is incorrect for COCO
        val2017 (order $10^3$ images). Training uses subsets and caption pools as
        documented per notebook.
  \item \textbf{KG construction:} the PDF's limitations mention word-frequency graphs
        while later sections assume NER-driven traversal. The aligned story is:
        \textbf{spaCy / scispaCy} entity extraction for general vs.\ medical domains,
        with PMI-weighted co-occurrence --- not a manually curated clinical ontology.
  \item \textbf{Privacy vs.\ APIs:} local Whisper and containerisation support a
        privacy narrative, but \textbf{Groq} (Llama inference) and \textbf{gTTS} are
        network services unless replaced with on-prem models. The proposal must not
        claim fully offline operation without listing exact local substitutes.
\end{itemize}

\section{Problem statement}

Voice-first assistants fail in high-stakes settings when they (i) hallucinate without
verifiable evidence, (ii) ignore non-text modalities, and (iii) cannot swap domain
knowledge without full retraining. MIVA-KNIGHT addresses (i) with retrieval-grounded
generation and confidence tiers, (ii) with audio and image encoders mapped into a
shared metric space, and (iii) with hot-swappable indices and domain configs while
keeping shared projection weights fixed per the capstone axioms.

\section{Research question (refined)}

\emph{How can dense multimodal embeddings, hybrid retrieval (vector index + corpus
graph signals), and a constrained LLM be composed so that answers remain tied to
retrieved evidence, while emotion and sentiment modules supply auxiliary context for
response style?}

\section{System architecture (as built)}

\subsection{Modality encoders and shared space}
Text (BERT-family), images (ResNet features + image projection), and audio (Wav2Vec
features + audio projection) map to unit-norm vectors in $\mathbb{S}^{511}$.
Domain transfer for radiology emphasises retraining or adapting the \emph{image
projection} head while freezing most backbone parameters, as detailed in Month~4 ROCO
documentation.

\subsection{Runtime pipeline (assistant loop)}
\begin{enumerate}[label=\textbf{Step \arabic*.}]
  \item Ingest text and/or speech (speech may be transcribed with Whisper-class STT).
  \item Encode query into the shared embedding space.
  \item Retrieve a shortlist with FAISS; re-rank with knowledge-graph signals derived
        from entity co-occurrence in the active corpus.
  \item Apply domain-specific confidence thresholds (traffic-light tiers).
  \item Generate with Llama-class LLM using only top-$k$ retrieved passages.
  \item Optional TTS; optional display of source snippets for auditability.
\end{enumerate}

\subsection{Offline training tracks (pipelines A--E)}
\begin{center}
\small
\begin{tabular}{@{}llp{6.5cm}@{}}
\toprule
\textbf{Track} & \textbf{Data} & \textbf{Purpose} \\
\midrule
A & COCO & General image--text alignment; general index + KG. \\
B & ROCO & Medical image--caption alignment; medical index + KG; surgical visual adapter. \\
C & RAVDESS & Audio emotion prototypes; \textbf{not} jointly trained with COCO. \\
D & CREMA-D & Audio--video emotion; InfoNCE / SupCon phases; warm-start from C. \\
E & CMU-MOSI & Three-stream sentiment fusion (text + audio + video frame); distinct from RAG retrieval fusion. \\
\bottomrule
\end{tabular}
\end{center}

\subsection{Notebooks: scope}
The \texttt{MIVA\_KNIGHT\_Challenges\_Implementation*.ipynb} notebooks demonstrate five
multimodal ML challenges (representation, translation proxies, fusion heads, co-learning
contrasts, alignment) on CMU-MOSI and CREMA caches. The
\texttt{MIVA\_KNIGHT\_RL\_Implementation\_*.ipynb} notebooks implement DQN-family
\emph{toy} control loops with multimodal state encoders as pedagogical complements to
the RL thesis chapter; they are not claims of deployed robotics.

\section{Objectives (measurable and honest)}

\begin{itemize}[nosep]
  \item \textbf{Retrieval quality:} report $\mathrm{P@}k$, MRR on held-out caption
        retrieval for COCO and ROCO as actually computed (including ablations for KG
        contribution where implemented).
  \item \textbf{Emotion / sentiment:} report per-class accuracy or MAE on RAVDESS,
        CREMA-D, and CMU-MOSI with the documented train/validation protocol (e.g.\
        video-level splits for MOSI).
  \item \textbf{Grounding:} qualitative and quantitative checks that answers cite
        retrieved spans; abstain rate under red-tier thresholds.
  \item \textbf{Engineering:} reproducible artifact paths (indices, checkpoints,
        caches), container images or environment specs, and latency of domain swap
        (index reload) where measured.
\end{itemize}

\noindent
\textbf{Explicit non-goals for the capstone prototype:} certifying HIPAA/GDPR
compliance, hitting unpublished third-party benchmark SOTA, or claiming real-time
streaming audio throughout without measured latency budgets.

\section{Evaluation plan}

\begin{enumerate}[nosep]
  \item Retrieval benchmarks on fixed validation draws from COCO and ROCO.
  \item Emotion benchmarks on RAVDESS / CREMA-D with confusion matrices.
  \item MOSI sentiment metrics separating regression (MAE) from classification (Acc-2),
        with notes on pseudo-label or label-noise issues when applicable.
  \item Optional \textbf{future work}: HaluBench-style hallucination probes and multi-hop
        QA sets once a stable generator--retriever interface is frozen.
\end{enumerate}

\section{Risks and mitigations}

\begin{description}[style=nextline,font=\normalfont\bfseries]
  \item[Modality interference.] Avoid joint contrastive training without true multimodal
        pairs; keep affect training on emotion datasets isolated from caption noise.
  \item[Cloud dependence.] Document every external API; provide a ``local-only''
        configuration table (STT, LLM, TTS substitutes).
  \item[Over-claiming fusion.] Reserve ``cross-attention fusion'' language for
        Pipeline~E (or explicit fusion heads), not for the entire MIVA-KNIGHT stack.
  \item[Medical safety.] ROCO is an \emph{IR} corpus; outputs are not diagnostic
        devices. Add institutional review and clinician-in-the-loop for any clinical use.
\end{description}

\section{Deliverables}

\begin{itemize}[nosep]
  \item LaTeX technical reports and thesis chapters as in the repository.
  \item Trained checkpoints, FAISS indices, KG serialisations, embedding caches.
  \item Notebooks with README paths for Google Drive / Colab assets.
  \item Demo script: voice or text query $\rightarrow$ retrieval trace $\rightarrow$
        grounded answer $\rightarrow$ optional emotion-conditioned phrasing.
\end{itemize}

\section{Conclusion}

MIVA-KNIGHT \emph{is} a multimodal system in the engineering sense: multiple encoders,
shared geometry for retrieval, hybrid sparse--dense evidence selection, and LLM
synthesis with guardrails. It is \emph{not} a single unified transformer trained
end-to-end on all datasets simultaneously. This aligned proposal matches that reality,
preserves the scientific contribution, and removes benchmark and dataset claims that
the repository does not substantiate.

\end{document}
